\subsection{筛}

定义 7. 筛是一个序列 \(\mathbf{C} = (\mathbf{C}_n, p_n)_{n \geqslant 0}\) ，使得对每个 \(n\) ，\(\mathbf{C}_n\) 是可数集，而 \(p_n\) 是 \(\mathbf{C}_{n+1}\) 到 \(\mathbf{C}_n\) 上的满射。

对每对整数 \(m, n\) （满足 \(0 \leqslant m \leqslant n\) ），令 \(p_{mn}\) 表示 \(C_m\) 到自身的恒等映射（若 \(m = n\) ），或映射 \(p_m \circ p_{m+1} \circ \cdots \circ p_{n-1}\) （它是 \(C_n\) 到 \(C_m\) 上的满射，若 \(m < n\) ）。显然 \(p_{mq} = p_{mn} \circ p_{ng}\) 对 \(m \leqslant n \leqslant q\) 成立，因此可以考虑 \(L(C)\) ，即族 \((C_n)\) 关于映射族 \((p_{mn})\) 的逆极限（集合论，第 III 章，§ 7）。若每个 \(C_n\) 赋以离散拓扑，则 \(L(C)\) 是拓扑空间的逆极限（第 I 章，§ 4，第 4 号）；作为这样的空间，\(L(C)\) 称为与筛 \(C\) 相伴的拓扑空间。\(L(C)\) 是拓扑积 \(\prod_n C_n\) 的闭子空间，由此立得 \(L(C)\) 是零维波兰空间（第 4 号）。

度量空间 X 的一个筛分由一个筛 \(\mathbf{C} = (\mathbf{C}_{n}, p_{n})\) 以及对每个整数 \(n \geqslant 0\) 的一个映射 \(\varphi_{n}\) 组成，后者把 \(C_{n}\) 映到 X 的直径 \(\leqslant 2^{-n}\) 的非空闭子集所成的集中，且满足：

\begin{enumerate}
\def\labelenumi{\alph{enumi})}
\item
  X 是诸集 \(\varphi_0(c)\) （当 \(c\) 跑遍 \(\mathbf{C}_0\) ）的并；
\item
  对每个 \(n\) 与每个 \(c \in \mathbf{C}_n\) ，\(\varphi_n(c)\) 是诸集 \(\varphi_{n+1}(c')\) （其中 \(c'\) 跑遍 \(\overline{\pmb{p}}_n(c)\) ）的并。
\end{enumerate}

筛分称为严格筛分，如果此外对每个 n ，诸集 \(\varphi_{n}(c)\) （当 c 跑遍 \(C_{n}\) ）两两不相交。

引理 3. 每个可数型度量空间 X 都具有一个筛分。若 X 还是零维的，则 X 具有严格筛分。

首先注意，若 Y 是可数型度量空间，\(\varepsilon\) 是实数 \textgreater0 ，则存在 Y 的由直径 \(\leqslant\varepsilon\) 的集组成的可数覆盖（§2，第 8 号，命题 13）。此外，若 Y 是零维的，则存在这样的、由既开又闭的集组成的覆盖 \((\mathbf{V}_{n})\) ；若 \(W_{n}\) 是 \(V_{n}\) 与 \(\bigcap_{k<n}(\mathbf{Y}-\mathbf{V}_{k})\) 的交，则可见诸 \(W_{n}\) 是闭的、直径 \(\leqslant\varepsilon\) 、两两不相交且覆盖 X。无论如何，覆盖中非空集的闭包构成 Y 的一个可数覆盖，其元素是直径 \(\leqslant\varepsilon\) 的非空闭集。

设 X 是可数型度量空间。设 \(C_{0}\) 是 X 的一个由直径 \(\leqslant I\) 的非空闭集组成的可数覆盖的指标集（若 X 是零维的，这些集两两不相交）；\(\varphi_{0}\) 是使每个指标 \(c \in C_{0}\) 对应于覆盖中相应集的映射。设我们已定义了诸 \(C_{i}\) 、诸 \(\varphi_{i}\) 与满射映射 \(p_{i}: C_{i+1} \to C_{i}\) （对 \(i \leqslant n\) ），使得对这些指标条件 b) 满足。若 \(c \in C_{n}\) ，则空间 \(\varphi_{n}(c)\) 有一个由直径 \(\leqslant I/2^{n+1}\) 的非空闭集组成的可数覆盖（若 X {[}从而 \(\varphi_{n}(c)\) {]}是零维的，这些集两两不相交）；若 \(\mathrm{I}(c)\) 表示此覆盖的指标集，我们取 \(c_{n+1}\) 为诸集 \(\mathrm{I}(c)\) （当 c 跑遍 \(C_{n}\) ）之和；对每个 \(c' \in C_{n+1}\) ，令 \(p_{n}(c')\) 表示元素 \(c \in C_{n}\) 中满足 \(c' \in \mathrm{I}(c)\) 者，并令 \(\varphi_{n+1}(c')\) 表示指标为 \(c'\) 的集（在所考虑的 \(\varphi_{n}(c)\) 的覆盖中）。显然我们这样归纳地定义了 X 的一个筛分，且若 X 是零维的，此筛分是严格的；由此得引理 3 。

现在设 X 是可数型完备度量空间，考虑 X 的由一个筛 C 与诸映射 \(\varphi_{n}\) 给出的一个筛分。若 \(\gamma = (c_{n})\) 是与 C 相伴的空间 \(\mathrm{L}(\mathrm{C})\) 的一点，则序列 \((\varphi_{n}(c_{n}))\) 是 X 中闭集的递减序列，其直径趋于 0 ；因此这序列的集的交由单独一点组成，我们记之为 \(f(\gamma)\) 。这样我们就定义了一个映射 \(f: \mathrm{L}(\mathrm{C}) \to \mathrm{X}\) 。若 \(\gamma, \gamma'\) 是 \(\mathrm{L}(\mathrm{C})\) 的两点，它们对 \(i \leqslant n\) 有相同的第 i 坐标，则显然 \(f(\gamma)\) 与 \(f(\gamma')\) 之间的距离 \(\leqslant 1/2^{n}\) ，从而 f 是连续的（由 \(\mathrm{L}(\mathrm{C})\) 的拓扑的定义）。对每个 \(x \in X\) ，由筛分的定义可用对 n 的归纳定义一个序列 \(\gamma = (c_{n})\) ，使得 \(x \in \varphi_{n}(c_{n})\) 对每个 \(n \geqslant 0\) 成立，且 \(c_{n} = p_{n}(c_{n+1})\) ；因此 \(x = f(\gamma)\) ，故 f 是满射。此外，若筛分是严格的，则序列 \(\gamma = (c_{n})\) （它满足 \(x = f(\gamma)\) ）是唯一的，故在此情形 f 是双射。f 称为所考虑筛分诱导的映射。

命题 13. 若 X 是任意卢津（相应地，苏斯林）空间，则存在一个筛 C 与 L(C) 到 X 上的连续双射（相应地，满射）。

参照卢津空间的定义（第 4 号，定义 6）{[}相应地，苏斯林空间的定义（第 2 号，定义 2）{]}，我们归结到如下情形：

X 是零维波兰空间（相应地，波兰空间），而前面的论证随即证明此结果。

